\section{Mechanization}\label{sec:mech}

The decision procedure of Section~\ref{sec:decide}, and the criterion of
Lemma~\ref{lem:dense} on which validity rests, are proved in Coq~8.19, in
\evalCoqLines{} lines. No goal is admitted, and \texttt{Print Assumptions}
reports every theorem below closed under the global context, so none of them
depends on an axiom.

\begin{center}
\begin{tabular}{@{}lll@{}}
\toprule
Paper & Coq & File \\
\midrule
Lemma~\ref{lem:dense} (dense bijections) & \texttt{dense\_iff\_running} & \texttt{Dense.v} \\
Lemma~\ref{lem:separable} (separability) & \texttt{separable\_iff} & \texttt{Separable.v} \\
Lemma~\ref{lem:floor} (components as floors) & \texttt{dgsum\_fsum} & \texttt{Chain.v} \\
Theorem~\ref{thm:decide} (linear recognition) & \texttt{scan\_iff}, \texttt{scan\_sound\_layout} & \texttt{Complete.v}, \texttt{Shape.v} \\
Section~\ref{sec:decide} as a whole & \texttt{decide\_iff} & \texttt{Decide.v} \\
\bottomrule
\end{tabular}
\end{center}

Lemma~\ref{lem:dense}, Lemma~\ref{lem:separable} and Theorem~\ref{thm:decide}
are each proved as an \emph{if and only if}. The last row composes them into a
statement about the checks the implementation runs. Take $k$ to be the value
at the origin and $g_i(v) = f(v\,e_i) - k$, check \eqref{eq:sep} in one pass
over the box, and run the scan of Theorem~\ref{thm:decide} on each $g_i$:
\begin{lstlisting}
Theorem decide_iff (S : list nat) (f : list nat -> Z) :
  sizes_pos S ->
  (IsRefined S f <-> Separable S f /\ AllAccept S (axis_gs S f)).
\end{lstlisting}
where \texttt{IsRefined S f} says that $f$ is the index function of a layout
over a refinement of $S$, plus a constant. The two directions do different
work. Soundness is what stops a wrong address expression from being emitted.
Completeness is what makes the procedure a decision rather than a heuristic:
no simplification is missed.

\paragraph{Where the proofs differ from the text.} The scan is phrased with
its residue recomputed from the widths recorded so far, rather than by
subtracting from a copy $\rho$ of $h$. By the invariant in the proof of
Theorem~\ref{thm:decide} the two are the same number at every step. The
forward direction of Lemma~\ref{lem:dense} avoids the counting step of its
proof in favour of an alignment argument: disjoint blocks of $n$ consecutive
positions that cover $[0,M)$ must be the aligned ones, with the covering given
by surjectivity and the disjointness by injectivity.

\paragraph{The proved scan is the one that runs.} The scan and the
construction of its shape are extracted from Coq to OCaml
(\texttt{coq/extract.sh}) and run by the test suite beside the
implementation's scan on every map of the enumeration in
Section~\ref{sec:eval}. On all \evalExtractedMaps{} maps they agree on the
verdict and, on the \evalExtractedStrided{} they accept, digit for digit. The
extraction remaps no datatype onto machine integers, so it assumes nothing the
proofs do not establish.

\paragraph{Not mechanized.} The minimality of the coarsest chain, the last
claim of Theorem~\ref{thm:decide}, which the tests check against a brute-force
search over divisor chains. The identity $s_j = g(w_j)$ of
Lemma~\ref{lem:floor}. The regrouping of a nested shape into the flat shape of
its leaves, and the sorting of modes by stride in Lemma~\ref{lem:dense}, the
proofs being stated for flat shapes with modes already sorted. And the
operations of Section~\ref{sec:ops}, which the tests check, with their laws,
against enumeration.
